\subsection{1. Đặt vấn đề}

\indent\indent Trong kỷ nguyên Web3, các giá trị của mạng xã hội phi tập trung (SocialFi), tiêu biểu là Lens Protocol\footnote{\url{https://lens.xyz/}}, ngày càng được khẳng định nhờ cơ chế trao hoàn toàn quyền sở hữu dữ liệu và khả năng tương tác sinh lời cho người dùng, góp phần bảo vệ lợi ích cộng đồng và mở ra cơ hội phát triển cho quản trị phi tập trung. Tuy nhiên, tính đặc thù này cũng chính đặc thù này khiến Web3 trở thành mục tiêu trọng điểm của tấn công Sybil, hình thức phá hoại hệ thống bằng cách sử dụng một số lượng lớn danh tính ẩn danh để tạo lập tầm ảnh hưởng \cite{al2017sybil}. Trong SocialFi, cuộc tấn công này biểu hiện dưới các hình thức phức tạp và tinh vi hơn khi kẻ tấn công sử dụng số lượng lớn tài khoản không chỉ để phát tán thông tin sai lệch mà còn trực tiếp chiếm đoạt tài nguyên từ các sự kiện Airdrop, lạm dụng cơ chế thưởng tương tác và làm sai lệch kết quả bỏ phiếu trong các tổ chức tự trị phi tập trung \cite{liu2025detecting}. Sự hiện diện của Sybil, gây thiệt hại nghiêm trọng về nguyên tắc phi tập trung và sự công bằng.\vspace{0.3cm}

Đặc điểm của dữ liệu trên mạng xã hội phi tập trung là tính đa quan hệ. Mỗi thực thể người dùng thường tồn tại đồng thời dưới nhiều tầng dữ liệu: thông tin hồ sơ, hành vi tương tác xã hội, và đặc biệt là tính sở hữu định danh (đồng sở hữu). Bên cạnh đó, các mạng lưới Sybil ngày càng tinh vi trong việc ngụy trang thông tin hồ sơ cá nhân để mô phỏng người dùng thật, làm tăng độ khó trong việc phân lớp nếu chỉ dựa vào một phương thức dữ liệu đơn lẻ hoặc các cấu trúc đồ thị truyền thống. Những hạn chế này khiến các phương pháp học máy cơ bản như máy học truyền thống hay mạng nơ-ron đa tầng bỏ sót các mối liên hệ ẩn và các ràng buộc cấu trúc quan trọng trong không gian đồ thị.\vspace{0.3cm}

Trước những thách thức trên, đề tài đề xuất hệ thống \textit{SybilSignal}, ứng dụng mô hình trí tuệ đồ thị nhận thức ràng buộc (Constraint-Aware) với quy trình ba giai đoạn. Đầu tiên, bộ mã hóa đồ thị (Graph Autoencoder, viết tắt là GAE) \cite{kipf2016variational} kết hợp cùng mạng nơ-ron đồ thị chú ý cải tiến (Graph Attention Network v2, viết tắt là GATv2) \cite{brody2021attentive} được sử dụng để trích xuất đặc trưng cấu trúc và học không giám sát biểu diễn nút. Tiếp đó, thuật toán gom cụm K-Means cùng hệ thống luật chuyên gia để tính điểm và gán nhãn giả lập dựa trên mức độ rủi ro, giải quyết triệt để sự thiếu hụt nhãn thực tế; đồng thời, tính giải thích được đảm bảo thông qua việc lưu vết lý do vi phạm và tận dụng lý do gán nhãn của láng giềng trên đồ thị tham chiếu để minh chứng cho kết quả dự đoán. Điểm đột phá của nghiên cứu nằm ở việc thực hiện đánh giá loại trừ, so sánh các kiến trúc mạng nơ-ron đồ thị (Graph Neural Network, viết tắt là GNN) \cite{ivanov2021boosted} cơ sở với mô hình tổ hợp giữa GNN và các bộ phân lớp học máy truyền thống. Toàn bộ quy trình được kiểm chứng trên ba phạm vi dữ liệu thực tế (tuần, tháng, quý). Sự kết hợp này giúp tối ưu hóa hiệu năng phát hiện Sybil, góp phần bảo vệ an toàn và bền vững cho hệ sinh thái mạng xã hội phi tập trung.v 